% =====================================================================
%  EurekaZero -- Section 5: The Price of Play
%  File: sections/05_price.tex   (compiles under the shared preamble)
% =====================================================================
\section{The Price of Play}\label{sec:price}
% Reader question: Can anyone afford to play it?
% One-sentence answer: Yes; cost per win is cost per attempt divided by the
%   win rate, so wasted moves set the bill, not token prices.
% Core points advanced: 3 (search = 1) and 2 (search length is the
%   reciprocal of taste); supports 1 (cheap replay is what invention lacks).
% Exhibits: Equation 2 (eq:bill), Table 5 (tab:price), Figure 4 (fig:price), Box 3.
% Numerical scenarios are retained from the supplied Section 5 source drafts.
% A reproducible calculation artifact remains pending.

What would it cost to play the innovation game? Under the planning
assumptions below, generation and updates for a 30B-class open-weight writer
cost about \$12,000 in GPU time. This is a partial compute projection, not
an observed price for a complete training run. Encoding, input tokens,
repairs and referee execution must be added to the ledger. Cheap
replay is the one property of Go that invention lacks (\cref{sec:game}). Here
we separate planning assumptions from the measurements still required in
\cref{sec:results}.

\subsection{The bill in one line}\label{sub:price-bill}

For repeated attempts with a fixed win probability, expected cost per win
is mean cost per attempt, repairs included, divided by that probability:
\begin{equation}
  C_{\text{win}} = \frac{C_{\text{attempt}} + \mathbb{E}[r]\,C_{\text{repair}}}{q}.
  \label{eq:bill}
\end{equation}
Here $C_{\text{attempt}}$ is the cost of one try, including the referee's
fit; $C_{\text{repair}}$ is the cost of one repair turn, and $\mathbb{E}[r]$
the mean number of repairs per attempt. The $q$ is that of \cref{eq:taste},
the probability that one try beats the incumbent: what sets the number of
tries also sets the price. It factors as $q = q_{\text{run}}\,q_{\text{beat}}$,
the share of moves that run times the share of those that win. A move that
runs and loses still costs money, so we divide by $q$, not by the survival
rate. Generating a token costs about 2~FLOPs per active parameter, and
training on one about~6 \citep{kaplan2020scaling}. We give self-hosted
projections and dated API-price examples (\cref{sub:price-run}).

The bill has two ledgers. Training pays for every attempted candidate,
including failures; deployment pays for encoding, one proposal and its
fitting procedure, with no repair. Fixed-weight agent loops instead buy
evaluation during task-specific search; a trained proposer may amortize
some of that work across later tasks. Whether the saving repays training
requires a measured crossover, still pending in \cref{sec:results}.

\subsection{Wasted moves}\label{sub:price-waste}

An attempt can fail before or after it becomes a valid competitor.
A stillborn move never competes: it fails
to parse or type-check, crashes, times out or returns degenerate output. An
also-ran runs and loses. The planned funnel in \cref{fig:price}a traces
those stages and the compute spent before each failed attempt stops. If 90\% of attempts fail to improve the incumbent, the independent-try
model requires ten attempts per improvement in expectation. A typed program language
removes many stillborn moves by construction (\cref{sec:game}), and those
that remain usually fail before the costly fit. Whether also-rans dominate is a pilot question. Learning, a changed search
procedure or a changed opponent can affect the measured win rate; isolate
the effect attributed to the writer.

\begin{figure}[t]
  \centering
  \includegraphics[width=\linewidth]{figures/fig4_price.pdf}
  \caption{\textbf{Pilot cost measurements pending.} Planned panels:
  (a)~the share of development proposals that run, score, beat the incumbent
  and survive re-verification, with the compute spent at each stage;
  (b)~cost per verified win by writer size and reasoning effort.
  \TBD{Supply pilot data, seeds and family-level uncertainty; identify
  settings with no verified win without assigning a finite cost per win.}}
  \label{fig:price}
\end{figure}

\subsection{The correction tax and the reasoning dial}\label{sub:price-tax}

Small models are cheap per token and can be dear per win. Compute per token
falls about 190-fold from a dense 1.5T writer to an 8B one, so the small
model is cheaper per win unless its win rate per attempt is more than 190
times lower, repairs counted. That threshold flatters the small model:
repair turns re-read a growing context, and a writer that needs three times
the tokens per attempt breaks even at about 60, not 190.

Reasoning is the second dial. Reasoning tokens may cut stillborn moves, but
each is billed as output. \Cref{fig:price}b maps cost per win over writer
size and reasoning effort. A U-shaped curve is the hypothesis; the
measurement decides where, and whether, the minimum sits.

\subsection{Rewarding useful hallucination}\label{sub:price-hallucination}

A pretrained writer pulls toward the textbook answer, which is history's
taste. On a board whose incumbent is that answer, it is an also-ran by
definition. A \emph{useful hallucination} is a move the pretrained prior rates
unlikely and the referee confirms (Box~3). Four counterweights push the
writer toward such moves: a novelty bonus paid only inside the validity gate
(\cref{sec:game}), temperature and entropy schedules, a curriculum that
starts past known solutions, and a quality-diversity archive
\citep{lehman2011novelty,mouret2015mapelites}. \Cref{sec:method} states which
of these \EZ{} uses. Novelty also raises waste, so exploration gets a stated
share of games, \TBD{X\%}.

\begin{explainer}{Box 3: Why a referee makes hallucination safe}
  \label{box:hallucination}
  In text, a hallucination is an error nobody checks. In a game, an
  unfamiliar move is a conjecture, and the referee checks every one.
  Unfamiliar and verified is a discovery; unfamiliar and unchecked is
  confabulation. In March 2016, AlphaGo's own estimate put the chance that a
  human would play its Move~37 against Lee Sedol at about one in ten
  thousand \citep{metz2016twomoves}, and AlphaGo won the game. A referee does
  not make strange moves good. It makes them cheap to try, and it is only as
  trustworthy as what it measures (\cref{sec:method}).
\end{explainer}

\subsection{How many games?}\label{sub:price-games}

Three quantities differ by orders of magnitude.
\textbf{Tries to find one improvement:} $1/q$ on average (\cref{eq:taste}).
If tries were independent, $q = 1\%$ would need 69 tries for even odds and
299 for 95\% confidence; tries from one writer are correlated, so these are
guides, not guarantees.
\textbf{Experience to train a transferable player:} AlphaGo Zero generated
4.9 million self-play games in about three days \citep{silver2017mastering},
and AlphaZero played 44 million games of chess in about nine hours
\citep{silver2017chess,silver2018general}. Our planning budget of 10,000
problems, 100 rounds and 8 attempts gives 8 million attempts, the same order.
It counts work, not learning: problem families matter more, and pilot
learning curves decide, read only within their measured range.
\textbf{Precision to plan the test:} illustratively, 150 independent
units at a win rate near 60\% give an approximate 95\% interval of
$\pm 8$ points. The actual study must compute uncertainty over independent
families, rather than treating every dataset as an independent observation.

\subsection{A run, priced}\label{sub:price-run}

\Cref{tab:price} prices an 8-million-attempt run for six writers: about
\$12,000 with a 30B-class writer and about \$570,000 with a dense 1.5T one.
The dense 1.5T rollouts take $1.2\times10^{23}$~FLOPs, about 1\% of a $10^{25}$-FLOP
pretraining run, a scale more than 30 public models had reached by mid-2025
\citep{rahman2025over30,epoch2024data}. Mixture-of-experts (MoE) writers pay
compute on their active parameters, such as the 32B of 1T in Kimi~K2 \citep{kimiteam2025k2}
or the 37B of 671B in DeepSeek-V3 \citep{deepseekv3}. They need memory for
all of them: 1T one-byte weights fill at least 13 H100s.

\begin{table}[t]
  \centering
  \caption{\textbf{Projected generation and update costs for a fixed
  attempt budget.} These totals exclude other components of a full run;
  the number of verified wins remains to be measured. Projections before pilots; [low-cost, high-cost] in
  brackets. Base: 8 million attempts of 5,000 generated tokens; \$2.92 per
  H100-hour on 29 October 2026 \citep{ornn2026ocpi} [\$2.50, \$4.17;
  \citealp{thunder2026gpu}]; 989.5 dense BF16 TFLOP/s per GPU
  \citep{nvidia2026h100}; generation at 20\% of peak [30\%, 10\%], or 10\%
  [20\%, 7\%] for memory-bound mixture-of-experts (MoE) writers; updates on
  10\% of tokens at 40\% of peak [50\%, 30\%]. The 1.5T rows are scenarios.
  Excludes input tokens, repairs and the referee. \$k: thousands of dollars;
  assumptions and numerical scenarios from the supplied cost drafts.
  \TBD{Supply a reproducible calculation artifact before publication.}}
  \label{tab:price}
  \small
  \begin{tabular}{@{}l l S[table-format=2.2] S[table-format=2.2] l l@{}}
    \toprule
    Program writer & {Active} & {GPU-h per} & {\$ per} & {Rollouts, 8M} & {Updates,} \\
                   & {params} & {1,000 attempts} & {1,000 attempts} & {attempts (\$k)} & {full run (\$k)} \\
    \midrule
    Small dense & 8B & 0.11 & 0.33 & 2.6 [1.5, 7.5] & 0.4 [0.3, 0.7] \\
    Mid dense & 32B & 0.45 & 1.31 & 10.5 [6.0, 30.0] & 1.6 [1.1, 3.0] \\
    Large dense & 70B & 0.98 & 2.87 & 23.0 [13.1, 65.6] & 3.4 [2.4, 6.6] \\
    1T MoE & 32B & 0.90 & 2.62 & 21.0 [9.0, 42.8] & 1.6 [1.1, 3.0] \\
    1.5T MoE (scenario) & 50B & 1.40 & 4.10 & 32.8 [14.0, 66.9] & 2.5 [1.7, 4.7] \\
    1.5T dense (scenario) & 1.5T & 21.05 & 61 & 492 [281, 1405] & 73.8 [50.5, 140] \\
    \bottomrule
  \end{tabular}
\end{table}

The referee runs on CPUs, at 10 to 100 CPU-seconds per move and
\$0.015--\$0.045 per vCPU-hour \citep{doit2026c7i}. For a full run that is
\$330--\$990 at 10 seconds and \$3,300--\$9,900 at 100, enough to exceed an
8B writer's whole bill. List prices sit above our floor. On 30 October 2026,
Together AI listed \$0.60 per million output tokens for gpt-oss-120B (5.1B
active) and \$0.14 for an 8B Llama model
\citep{together2026pricing,openai2025gptoss}: about 7 and 2 times our floor
for the same model. The running example's single game cost \TBD{\$c}:
\TBD{t} generated tokens and \TBD{s} CPU-seconds of referee time.

\begin{keyidea}
  Cost per verified improvement depends on both the cost of an attempt
  and the probability that it wins.
\end{keyidea}

The table prices a fixed number of attempts, not a fixed number of wins.
Converting those totals to cost per win requires a measured win rate. Under
\cref{eq:bill}'s fixed-probability accounting, doubling $q$ halves the
expected cost per win while leaving the per-attempt bill unchanged.

The scaling scenarios in this section concern 8B--1.5T language-model
writers producing 5,000 tokens per attempt; they do not price the reference
checkpoint in \cref{sec:method}. That checkpoint uses a compact network with 27.8 million parameters and at most 128 grammar-constrained tokens, with no repair
in the proposed single-candidate deployment protocol. Its measured costs
must be reported separately from the scaling projections.

The method's job is to improve proposal quality; whether it does so must be
measured under the stated deployment budget.
